%! Rational Functions and Zeros — Unit 3, Lesson 3.2 — Guided Notes (Answer Key)
\documentclass[10pt]{article}
\usepackage{precalculus-article}
\usepackage{precalculus-key}   % pulls in -boxes; do NOT also load -boxes
\newcommand{\TallMath}[1]{$\displaystyle #1\rule[-1.4em]{0pt}{3.2em}$}

\begin{document}

\pageheader{Unit 3, Lesson 3.2}{Guided Notes --- Answer Key}
\namedateperiod

\vspace{0.10in}

\begin{objectivebox}
By the end of this lesson, I will be able to\ldots
\begin{enumerate}[itemsep=2pt]
  \item Find the real zeros of a rational function by setting the \ans{numerator} equal to zero
        and checking that the solutions are in the \ans{domain}. \hfill\textit{(LO 1.8.A)}
  \item Solve a rational inequality using a \ans{sign chart} with critical values from
        both the numerator and the denominator. \hfill\textit{(LO 1.8.A)}
\end{enumerate}
\end{objectivebox}

\vspace{0.08in}

\begin{vocabbox}
\termblanklong{Zero of a rational function}
\ansline{An input value $c$ in the domain where $r(c) = 0$; equivalently, where the numerator equals zero and the denominator does not.}
\termblanklong{Domain restriction}
\ansline{A value excluded from the domain because it makes the denominator zero; NOT a zero of the function.}
\termblanklong{Critical value}
\ansline{Any $x$-value where $r(x) = 0$ (numerator zero in domain) or $r(x)$ is undefined (denominator zero); used as interval boundaries in sign analysis.}
\termblanklong{Sign chart}
\ansline{A number-line diagram that marks critical values and records the sign ($+$ or $-$) of $r(x)$ on each resulting interval.}
\end{vocabbox}

\vspace{0.08in}

\begin{hookbox}
A company's profit per unit is modeled by $\displaystyle P(x) = \dfrac{x^2 - 9}{x - 1}$
where $x > 0$, $x \ne 1$.

At what production level(s) is the profit per unit equal to zero? When is it positive?
Why does $x = 1$ need to be excluded?

\vspace{0.05in}
\ansline{Set numerator $= 0$: $x^2 - 9 = 0 \Rightarrow x = 3$ (in domain). Profit is zero at $x = 3$. Sign analysis gives positive profit on $(1, 3)$ would need checking --- $x = 1$ is excluded because $P(1)$ is undefined (denominator $= 0$).}
\end{hookbox}

\vspace{0.08in}

\begin{notesbox}{Zeros of Rational Functions (EK 1.8.A.1)}
\textbf{Rule:} For $r(x) = \dfrac{p(x)}{q(x)}$, we have $r(x) = 0$ when:
\begin{itemize}[itemsep=3pt]
  \item $p(x) = 0$ (numerator equals \ans{zero}), AND
  \item $q(x) \ne 0$ ($x$ is in the \ans{domain} of $r$).
\end{itemize}

\medskip
\textbf{Worked Example.} Find all zeros of $r(x) = \dfrac{x^2 - 4}{x - 3}$.

\medskip
\renewcommand{\arraystretch}{1.5}
\begin{tabular}{@{} r l @{\qquad} l @{}}
\textbf{Step 1} & Find domain restrictions. & $q(x) = 0 \Rightarrow$ \ans{$x - 3 = 0 \Rightarrow x = 3$} \\
                &                           & Domain: all reals except $x =$ \ans{$3$} \\[4pt]
\textbf{Step 2} & Set numerator $= 0$.      & $x^2 - 4 = 0 \Rightarrow (x - \ans{2})(x + \ans{2}) = 0$ \\
                &                           & Solutions: $x =$ \ans{$2$} and $x =$ \ans{$-2$} \\[4pt]
\textbf{Step 3} & Check domain.             & Is $x = 2$ in the domain? \ans{Yes} \\
                &                           & Is $x = -2$ in the domain? \ans{Yes} \\[4pt]
\textbf{Step 4} & State zeros.              & Zeros of $r$: $x =$ \ans{$2$ and $x = -2$} \\
\end{tabular}
\end{notesbox}

\vspace{0.08in}

\begin{notesbox}{Rational Inequalities --- Sign Chart Method (EK 1.8.A.2)}
\textbf{Key idea:} The sign of $r(x)$ can only change at a \ans{numerator zero (in domain)} or a
\ans{denominator zero (domain restriction)}. These are called \textbf{critical values}.

\medskip
\textbf{Solve:} $\dfrac{x - 1}{x + 2} \geq 0$

\medskip
\renewcommand{\arraystretch}{1.5}
\begin{tabular}{@{} r l @{\qquad} l @{}}
\textbf{Step 1} & Find critical values. & Numerator zero: $x =$ \ans{$1$} \quad Denominator zero: $x =$ \ans{$-2$} \\[4pt]
\textbf{Step 2} & Mark on number line.  & Open circle at $x =$ \ans{$-2$} (excluded); closed at $x =$ \ans{$1$} (zero). \\[4pt]
\textbf{Step 3} & Test each interval.   & \\
\end{tabular}

\medskip
\begin{center}
\renewcommand{\arraystretch}{1.6}
\begin{tabular}{|c|c|c|c|}
\hline
\textbf{Interval} & $(-\infty,\,-2)$ & $(-2,\,1)$ & $(1,\,\infty)$ \\
\hline
Test value $x$    & \ans{$-3$}  & \ans{$0$}  & \ans{$2$} \\
\hline
Sign of $r(x)$    & \ans{$(+)$} & \ans{$(-)$} & \ans{$(+)$} \\
\hline
\end{tabular}
\end{center}

\medskip
\textbf{Step 4.} Write the solution (intervals where $r(x) \geq 0$):

\vspace{0.06in}
Solution: \ans{$(-\infty,\,-2)\cup[1,\,\infty)$}
\end{notesbox}

\end{document}
